% Pursuing Stacks (cotes 134-2 to 134-8): archive cote 134-7, site page 104 (PDF page 105), typescript page 546.
% typescript. Draft compared: ps6.tex:5231-5287.
% AI-transcribed by Claude Opus 5.5 (claude-opus-5-5, adaptive thinking, effort high), Message Batches API, batch ps4 (msgbatch_012Aoc1NuskATPsyBeKSYK7J), 2026-11-12; not yet checked by eye.
\scanpage{134-7/104}
\mspage{546}
reflections on schematization and on schematic homotopy types may make a
sense.

Definitely my suggestion in\inserted{the} last notes\hand{,} to work with non connected
``bundles''\hand{,} is'nt much more than the reflection of my unability so far to
make a breakthrough and get the\xout{right} ``left'' homotopy constructions
in terms of semisimplicial (connected) unipotent bundles alone. Besides
the serious drawback already pointed out at the end of section 130, ano-%
ther one occured to me - namely that with the suggested extension, one
is lo\unsure{\inserted{o}}sing track of the extra condition $X_0=e$ (or even $X_0=X_1=e$) on our
models, which is often remaining implicit in the notes, and which how-%
ever is an important restriction\hand{,} actually needed for the kind of things
we want to do. For one thing, when we drop it, and even when we other-%
wise restrict to components $X_n$ which are standard affine spaces $E^{d_n}_k$ say,
there is no hope of getting a Lie-type description for the homotopy shea-%
ves $\hand{\pi}_i(X_*)$ - for instance, the Lie-type invariants we get \xout{from} by
using different sections of $X_0$ over $S=\mathrm{Spec}(k)$ are by no means related
(as the\xout{y} \inserted{$\pi_i$} should) by a transitive system of canonical isomorphisms; this
is easily seen already when taking the ``trivial'' semisimplicial bundle
described by $X_n=X_0$ for all n (the constant functor with value $X_0$ from
$\hand{\Delta}^0$ to $U(k)$)! Of course, even when allowing connected components for
the ``bundles'' $X_n$ , we may formally still throw in the condition $X_0=e$
- but this is cheating and no use at all, if we remember that the main
motivation for allowing connected components was in order to be able (in
suitable cases) to perform \xout{in a standard way} an integration operation of
schematic homotopy types. But when following the standard construction,
for the resulting $X_*$ the component $X_0$ as well as all others will have
lots of connected components, and hence the condition $X_0=e$ will not
hold true. The same remark applies also to the ``constant'' semisimplicial
bundles (namely with \xout{constant} components which are ``constant'' schemes
over S ) defined by the standard simplices $\hand{\Delta}_n$ . To sum up, while it \uline{is}
certainly quite useful to view semisimplicial ``bundles'', used for descri-%
bing the ``schematic homotopy types'' we are after, as particular cases
of \inserted{\deleted{\illegible{1}} more general} semisimplicial sheaves on the fpqc site of $S=\mathrm{Spec}(k)$, we will probably
have to be quite careful in keeping the ``bundles'' we are working with
restricted enough, and not confuse our models for schematic homotopy
types (allowing for a nice description of the basic $L\pi_{.}$ and $LH_{.}$ invari-%
ants) with more general sem\hand{i}simplicial sheaves which may enter the picture
in various ways.

\margin{104}
